%% kapa: exam=1 round=35 year=2024 period=1 form=A subject=civil_law subject_ko=민법 seq=1 number=1
%% source: Q-Net 감정평가사 기출문제 (cst00309) artlSeq=5213229 file=raw/1008/5213229/1교시 문제지 원본.pdf
%% conversion: kordoc 4.24.51 -> tools/gen1.py (auto); review=proofread-sonnet; flags=none
\begin{question}{1}
민법의 법원(法源)에 관한 설명으로 옳지 않은 것은? (다툼이 있으면 판례에 따름)
\begin{choices}
  \item 민사에 관한 헌법재판소의 결정은 민법의 법원이 될 수 있다.
  \item 사적자치가 인정되는 분야의 제정법이 주로 임의규정인 경우, 사실인 관습은 법률행위 해석기준이 될 수 있다.
  \item 법원(法院)은 판례변경을 통해 기존 관습법의 효력을 부정할 수 있다.
  \item 관습법은 사회 구성원의 법적 확신으로 성립된 것이므로 제정법과 배치되는 경우에는 관습법이 우선한다.
  \item 법원(法院)은 관습법에 관한 당사자의 주장이 없더라도 직권으로 그 존재를 확정할 수 있다.
\end{choices}
\end{question}
\answer{4}
